\begin{table*}[!t]
\caption{ZAC18与QMAP154上三种编译方法的总体结果}
\label{tab:main-results}
\centering
\PaperTableSetup
\begin{PaperRevisionBlock}
\begin{tabular*}{\textwidth}{@{\extracolsep{\fill}}llrrrr@{}}
\toprule
数据集 & 方法 & 模型保真度 & 重排批次 & \shortstack{重排时延\\(ms)} & \shortstack{完整结果\\/$N$}\\
\midrule
\multirow{3}{*}{\shortstack{ZAC18\\($N_{\rm ind}=\DefaultZACStrictN$)}}
& ZAC & \pgfmathprintnumber[fixed,precision=4,zerofill]{\DefaultZACMOneF} & \DefaultZACMOneB & \DefaultZACMOneT & \DefaultZACMOneV\\
& ICCAD/QMAP & \pgfmathprintnumber[fixed,precision=4,zerofill]{\DefaultZACMTwoF} & \DefaultZACMTwoB & \DefaultZACMTwoT & \DefaultZACMTwoV\\
& GA-LK & \textbf{\boldmath\pgfmathprintnumber[fixed,precision=4,zerofill]{\DefaultZACMFourF}} & \textbf{\DefaultZACMFourB} & \textbf{\DefaultZACMFourT} & \DefaultZACMFourV\\
\midrule
\multirow{3}{*}{\shortstack{QMAP154\\($N_{\rm ind}=\DefaultQMAPStrictN$)}}
& ZAC & \pgfmathprintnumber[fixed,precision=6,zerofill]{\DefaultQMAPMOneF} & \DefaultQMAPMOneB & \DefaultQMAPMOneT & \DefaultQMAPMOneV\\
& ICCAD/QMAP & \pgfmathprintnumber[fixed,precision=6,zerofill]{\DefaultQMAPMTwoF} & \DefaultQMAPMTwoB & \DefaultQMAPMTwoT & \DefaultQMAPMTwoV\\
& GA-LK & \textbf{\boldmath\pgfmathprintnumber[fixed,precision=6,zerofill]{\DefaultQMAPMFourF}} & \textbf{\DefaultQMAPMFourB} & \textbf{\DefaultQMAPMFourT} & \DefaultQMAPMFourV\\
\bottomrule
\end{tabular*}
\end{PaperRevisionBlock}
\PaperTableNote[\textwidth]{注：保真度取几何均值，重排批次和时延取算术均值；粗体标示同一数据集内的最优值。$N_{\rm ind}$为三种方法共同有效的不同电路数，QMAP154的\PaperRevision{\DefaultQMAPStrictFileN} 个共同有效文件去重后对应\PaperRevision{\DefaultQMAPStrictN} 个电路。完整结果数按全部输入文件统计。}
\end{table*}
